% Part: normal-modal-logic
% Chapter: filtrations
% Section: S5-decidable

\documentclass[../../../include/open-logic-section]{subfiles}

\begin{document}

\olfileid{nml}{fil}{dec}

\olsection{\Log{S5} 是可判定的}
有限模型性质为我们提供了一种简便方法，用以证明由模式定义的模态逻辑系统是\emph{可判定的}
（即，存在一个可计算过程，能够判定任意公式在该系统中是否可推导）。

\begin{thm}
  \Log{S5} 是可判定的。
\end{thm}

\begin{proof}
  给定 $!A$，并设 $!A$ 中出现的命题变元均在 $p_1$, \dots, $p_k$ 之中。
  由于对每个 $n$，具有 $n$ 个世界且对 $p_1$, \dots, $p_k$ 赋值的模型只有有限多个，
  因此我们可以通过某种系统的方式生成证明，来\emph{并行地}枚举 \Log{S5} 的所有定理；
  同时枚举所有包含 $1$, $2$, \dots 个世界的模型，
  并检查 $!A$ 是否在某个此类模型的某个世界上失效。
  最终，这两个并行进程之一会给出答案，因为根据
  \olref[com][fra]{thm:generaldet} 和 \olref[fmp]{cor:S5fmp}，
  要么 $!A$ 是可推导的，要么 $!A$ 在某个有限全通模型中失效。
\end{proof}

上述证明适用于 \Log{S5}，是因为全通模型的过滤自动保持全通性。
这对自反性和序列性同样成立，但对于其他性质则需要更多工作。

\end{document}